\documentclass[../../main.tex]{subfiles}
\begin{document}

\begin{center}
  {\fontsize{48pt}{58pt}\selectfont \textbf{1980}\normalsize\textbf{年度}}
\end{center}

\begin{table}[htb]
  \centering
  \caption{出題テーマと難易度}
  \begin{tabular}{|c|p{9cm}|c|}
    \hline
    \textbf{問題番号} & \textbf{テーマ} & \textbf{難易度} \\
    \hline
    第1問 & 正三角形内の交点と面積比 & \\
    \hline
    第2問 & 球への射影で得られる曲線の長さ & \\
    \hline
    第3問 & 行列の累乗と数列の値 & \\
    \hline
    第4問 & 媒介変数で定まる点と原点の距離 & \\
    \hline
    第5問 & 正四面体の辺上の移動と確率 & \\
    \hline
    第6問 & 双曲線と二等辺三角形の共通部分の面積 & \\
    \hline
  \end{tabular}
\end{table}

\begin{table}[htb]
  \centering
  \caption{解答}
  \renewcommand{\arraystretch}{1.4}
  \begin{tabular}{|c|c|p{9cm}|}
    \hline
    \textbf{問題番号} & \textbf{小問} & \textbf{答え} \\
    \hline
    第1問 & (1) & $BB'=\dfrac{x}{\sqrt{x^2-x+1}}$，$PB'=\dfrac{x^2}{\sqrt{x^2-x+1}}$ \\
    \cline{2-3}
    & (2) & $x=\dfrac{7-\sqrt{21}}{14}$ \\
    \hline
    第2問 & -- & 曲線の長さは$\left(\dfrac{1}{2}+\dfrac{2\sqrt{6}}{9}\right)\pi$ \\
    \hline
    第3問 & (1) & \\
    \cline{2-3}
    & (2) & \\
    \hline
    第4問 & -- & 距離の2乗の最大値は$\dfrac{k^2}{16}+2$．最小値は$0<k\le2$なら$\dfrac{k^2}{16}$，$k\ge2$なら$1-\dfrac{3}{2}\sqrt[3]{\dfrac{1}{2k^2}}$ \\
    \hline
    第5問 & (1) & $p_1(n)=p_2(n)=p_3(n)$（証明問題） \\
    \cline{2-3}
    & (2) & $p_0(n)=\dfrac{1}{4}+\dfrac{3}{4}\left(-\dfrac{1}{3}\right)^n$，$p_1(n)=\dfrac{1}{4}\left\{1-\left(-\dfrac{1}{3}\right)^n\right\}$ \\
    \hline
    第6問 & -- & $0<s\le1$なら$s$，$s\ge1$なら$s-\sqrt{s^2-s}+\log\left(\sqrt{s}+\sqrt{s-1}\right)$ \\
    \hline
  \end{tabular}
\end{table}

\paragraph{試験形態}
6問．試験時間は要確認．

\end{document}
